\chapter{Projectile Motion as a Two-Dimensional Model}

\section{Purpose and Physical Assumptions}

Projectile motion is one of the simplest situations in which a two-dimensional
motion becomes understandable by separating it into two one-dimensional motions.
The horizontal motion is uniform, while the vertical motion is uniformly
accelerated by gravity.

In this note we derive the standard formulas for an ideal projectile launched
from a point with initial speed \(v_0\) at an angle \(\theta\) above the
horizontal. The purpose is not only to list the formulas, but to show where they
come from and what assumptions make them valid.

We use the following idealized model.

\begin{itemize}
    \item The projectile is treated as a point particle.
    \item Air resistance is neglected.
    \item The gravitational acceleration is constant and directed downward.
    \item The Earth is treated locally as flat, so horizontal and vertical axes
    are meaningful over the distance considered.
    \item The launch point is chosen as the origin unless stated otherwise.
\end{itemize}

With the positive \(y\)-axis pointing upward, the acceleration vector is
\[
\mathbf{a}=(0,-g),
\]
where \(g>0\) is the magnitude of gravitational acceleration.

\section{Geometry of the Launch}

Let the initial velocity be
\[
\mathbf{v}_0=(v_{0x},v_{0y}).
\]
If the initial speed is \(v_0=\lVert \mathbf{v}_0\rVert\) and the launch angle is
\(\theta\), then the velocity components are
\[
v_{0x}=v_0\cos\theta,
\qquad
v_{0y}=v_0\sin\theta.
\]
This is simply the decomposition of a vector into its horizontal and vertical
components.

\begin{figure}[h]
\centering
\begin{center}
\begin{tikzpicture}[scale=0.95, >=Stealth]
    \draw[->] (-0.2,0) -- (8.2,0) node[right] {\(x\)};
    \draw[->] (0,-0.2) -- (0,3.8) node[above] {\(y\)};

    \draw[thick, domain=0:7.2, samples=80]
        plot (\x,{0.85*\x - 0.095*\x*\x});

    \fill (0,0) circle (1.6pt) node[below left] {launch};
    \fill (7.2,0) circle (1.6pt) node[below] {landing};
    \fill (4.47,1.90) circle (1.6pt) node[above] {highest point};

    \draw[->, thick] (0,0) -- (1.8,1.53) node[midway, above left] {\(\mathbf{v}_0\)};
    \draw[dashed] (0,0) -- (2.1,0);
    \draw (0.65,0) arc[start angle=0,end angle=40,radius=0.65];
    \node at (0.95,0.28) {\(\theta\)};

    \draw[dashed] (4.47,0) -- (4.47,1.90);
    \draw[<->] (4.75,0) -- (4.75,1.90) node[midway,right] {\(H\)};
    \draw[<->] (0,-0.45) -- (7.2,-0.45) node[midway,below] {\(R\)};

    \draw[->] (6.35,2.6) -- (6.35,1.75) node[midway,right] {\(\mathbf{g}\)};
\end{tikzpicture}
\end{center}

\caption{Ideal projectile motion without air resistance. The trajectory is a
parabola, the acceleration is vertical, and the horizontal and vertical motions
can be treated separately.}
\end{figure}

\section{Equations of Motion from Constant Acceleration}

The acceleration equations are
\[
\frac{d^2x}{dt^2}=0,
\qquad
\frac{d^2y}{dt^2}=-g.
\]
We derive the position functions by integrating with respect to time.

For the horizontal coordinate,
\[
\frac{d^2x}{dt^2}=0.
\]
Integrating once gives
\[
\frac{dx}{dt}=C_1.
\]
The constant \(C_1\) is the initial horizontal velocity, so
\[
\frac{dx}{dt}=v_{0x}=v_0\cos\theta.
\]
Integrating again gives
\[
x(t)=v_{0x}t+C_2.
\]
If \(x(0)=0\), then \(C_2=0\), hence
\[
\boxed{x(t)=v_0\cos\theta\, t.}
\]

For the vertical coordinate,
\[
\frac{d^2y}{dt^2}=-g.
\]
Integrating once gives
\[
\frac{dy}{dt}=-gt+C_3.
\]
The constant \(C_3\) is the initial vertical velocity, so
\[
\frac{dy}{dt}=v_{0y}-gt=v_0\sin\theta-gt.
\]
Integrating again gives
\[
y(t)=v_{0y}t-\frac{1}{2}gt^2+C_4.
\]
If \(y(0)=0\), then \(C_4=0\), hence
\[
\boxed{y(t)=v_0\sin\theta\, t-\frac{1}{2}gt^2.}
\]

Thus the velocity and position are
\[
\boxed{\mathbf{v}(t)=\bigl(v_0\cos\theta,\,v_0\sin\theta-gt\bigr)}
\]
and
\[
\boxed{\mathbf{r}(t)=\bigl(v_0\cos\theta\,t,\,v_0\sin\theta\,t-
\tfrac{1}{2}gt^2\bigr).}
\]

\begin{remarkbox}
The horizontal coordinate changes linearly in time because there is no
horizontal acceleration in the ideal model. The vertical coordinate contains a
quadratic term because the vertical acceleration is constant and nonzero.
\end{remarkbox}

\section{The Trajectory as a Parabola}

The parametric equations are
\[
x(t)=v_0\cos\theta\,t,
\qquad
y(t)=v_0\sin\theta\,t-\frac{1}{2}gt^2.
\]
Assume \(\cos\theta\neq 0\). From the horizontal equation,
\[
t=\frac{x}{v_0\cos\theta}.
\]
Substituting this into \(y(t)\) gives
\[
y=x\tan\theta-
\frac{g}{2v_0^2\cos^2\theta}x^2.
\]
Therefore the trajectory is
\[
\boxed{y(x)=x\tan\theta-
\frac{g}{2v_0^2\cos^2\theta}x^2.}
\]
This is a quadratic function of \(x\), so the ideal trajectory is a parabola.

\section{Time of Flight}

The projectile returns to the launch height when \(y(t)=0\). With launch and
landing at the same height,
\[
y(t)=v_0\sin\theta\,t-\frac{1}{2}gt^2=0.
\]
Factoring out \(t\),
\[
t\left(v_0\sin\theta-\frac{1}{2}gt\right)=0.
\]
There are two solutions. The first is \(t=0\), the launch moment. The second is
found from
\[
v_0\sin\theta-\frac{1}{2}gt=0.
\]
Hence the total time of flight is
\[
\boxed{T=\frac{2v_0\sin\theta}{g}.}
\]
This formula applies when the projectile lands at the same vertical level from
which it was launched.

\section{Maximum Height}

At the highest point the vertical velocity is zero:
\[
v_y(t)=v_0\sin\theta-gt=0.
\]
Therefore the time to reach the top is
\[
\boxed{t_H=\frac{v_0\sin\theta}{g}.}
\]
Substituting \(t_H\) into \(y(t)\),
\[
H=v_0\sin\theta\left(\frac{v_0\sin\theta}{g}\right)
-\frac{1}{2}g\left(\frac{v_0\sin\theta}{g}\right)^2.
\]
Thus
\[
H=\frac{v_0^2\sin^2\theta}{g}-\frac{v_0^2\sin^2\theta}{2g}
=\frac{v_0^2\sin^2\theta}{2g}.
\]
So the maximum height above the launch point is
\[
\boxed{H=\frac{v_0^2\sin^2\theta}{2g}.}
\]

\section{Horizontal Range}

The range \(R\) is the horizontal distance traveled during the total time of
flight. Since
\[
x(t)=v_0\cos\theta\,t,
\]
we have
\[
R=x(T)=v_0\cos\theta\,T.
\]
Using \(T=2v_0\sin\theta/g\),
\[
R=v_0\cos\theta\cdot\frac{2v_0\sin\theta}{g}
=\frac{2v_0^2\sin\theta\cos\theta}{g}.
\]
Since \(2\sin\theta\cos\theta=\sin(2\theta)\),
\[
\boxed{R=\frac{v_0^2\sin(2\theta)}{g}.}
\]

The range is maximal when \(\sin(2\theta)=1\), that is when
\[
2\theta=90^\circ,
\qquad
\theta=45^\circ.
\]
For launch and landing at the same height, the ideal range is therefore maximal
at \(45^\circ\).

\section{Velocity, Speed, and Direction During Flight}

The velocity components are
\[
v_x(t)=v_0\cos\theta,
\qquad
v_y(t)=v_0\sin\theta-gt.
\]
The speed is the magnitude of the velocity vector:
\[
\lVert \mathbf{v}(t)\rVert
=\sqrt{v_x(t)^2+v_y(t)^2}
=\sqrt{v_0^2\cos^2\theta+\bigl(v_0\sin\theta-gt\bigr)^2}.
\]
Thus
\[
\boxed{\lVert \mathbf{v}(t)\rVert
=\sqrt{v_0^2\cos^2\theta+\bigl(v_0\sin\theta-gt\bigr)^2}.}
\]

The instantaneous direction angle \(\alpha(t)\) of the velocity satisfies
\[
\tan\alpha(t)=\frac{v_y(t)}{v_x(t)}
=\frac{v_0\sin\theta-gt}{v_0\cos\theta},
\]
provided \(v_0\cos\theta\neq 0\).

At the highest point, \(v_y=0\), so the velocity is purely horizontal:
\[
\mathbf{v}(t_H)=(v_0\cos\theta,0).
\]

\section{Projectile Launched from an Initial Height}

If the projectile is launched from height \(y_0\), the position equations become
\[
x(t)=x_0+v_0\cos\theta\,t,
\]
and
\[
y(t)=y_0+v_0\sin\theta\,t-\frac{1}{2}gt^2.
\]
If the ground is at \(y=0\), the landing time is obtained from
\[
y_0+v_0\sin\theta\,t-\frac{1}{2}gt^2=0.
\]
Equivalently,
\[
gt^2-2v_0\sin\theta\,t-2y_0=0.
\]
The physically relevant positive solution is
\[
\boxed{T=\frac{v_0\sin\theta+\sqrt{v_0^2\sin^2\theta+2gy_0}}{g}.}
\]
The corresponding range, measured from \(x_0\), is
\[
\boxed{R=v_0\cos\theta\,T.}
\]

When \(y_0=0\), the formula reduces to
\[
T=\frac{v_0\sin\theta+|v_0\sin\theta|}{g}.
\]
For an upward launch with \(\sin\theta>0\), this becomes
\[
T=\frac{2v_0\sin\theta}{g},
\]
which agrees with the same-height result.

\section{Energy Check}

The same maximum height formula can be checked using conservation of mechanical
energy. At launch, the vertical part of the kinetic energy associated with
upward motion is
\[
\frac{1}{2}m(v_0\sin\theta)^2.
\]
At the highest point, this vertical kinetic energy has been converted into
gravitational potential energy \(mgH\). Thus
\[
\frac{1}{2}m v_0^2\sin^2\theta=mgH.
\]
Canceling \(m\) and solving for \(H\),
\[
H=\frac{v_0^2\sin^2\theta}{2g}.
\]
This agrees with the kinematic derivation.

\section{A Worked Symbolic Example}

Suppose a projectile is launched from ground level with speed \(v_0=20\,\mathrm{m/s}\)
at angle \(\theta=30^\circ\). Since
\[
\sin 30^\circ=\frac{1}{2},
\qquad
\cos 30^\circ=\frac{\sqrt{3}}{2},
\]
the initial components are
\[
v_{0x}=10\sqrt{3}\,\mathrm{m/s},
\qquad
v_{0y}=10\,\mathrm{m/s}.
\]
The time of flight is
\[
T=\frac{2v_{0y}}{g}=\frac{20}{g}.
\]
The maximum height is
\[
H=\frac{v_{0y}^2}{2g}=\frac{100}{2g}=\frac{50}{g}.
\]
The range is
\[
R=v_{0x}T=10\sqrt{3}\cdot\frac{20}{g}=\frac{200\sqrt{3}}{g}.
\]
If one later substitutes \(g\approx 9.81\,\mathrm{m/s^2}\), these symbolic
expressions give numerical predictions in SI units.

\section{Limitations of the Model}

The ideal projectile model is powerful because it is simple, but its assumptions
are restrictive. In real motion, air resistance may reduce the range, change the
optimal launch angle, and make the trajectory non-parabolic. For high speeds or
long distances, the curvature of the Earth and variation of \(g\) may also
matter. The formulas derived here are therefore best understood as exact results
inside an ideal model, and as approximations for real projectiles when the
assumptions are sufficiently accurate.

\begin{summarybox}
Projectile motion without air resistance separates into horizontal uniform
motion and vertical uniformly accelerated motion. From
\(\mathbf{a}=(0,-g)\), integration gives
\[
x(t)=v_0\cos\theta\,t,
\qquad
y(t)=v_0\sin\theta\,t-\frac{1}{2}gt^2.
\]
Eliminating time gives a parabolic trajectory. For launch and landing at the
same height,
\[
T=\frac{2v_0\sin\theta}{g},
\qquad
H=\frac{v_0^2\sin^2\theta}{2g},
\qquad
R=\frac{v_0^2\sin(2\theta)}{g}.
\]
\end{summarybox}
